% =========================================================
\section{Integral Compressible Boundary-Layer Model for Quick Screening}
\label{sec:integral_boundary_layer}
% =========================================================

The internal flow of a rocket nozzle is characterized by high velocities,
strong streamwise acceleration, and large variations in pressure, temperature,
density, and viscosity. Furthermore, the generally high Reynolds numbers make
the near-wall flow predominantly turbulent. Consequently, an incompressible
flat-plate correlation is not sufficient to represent the boundary-layer
development inside the nozzle.

Resolving the complete viscous flow would require a two-dimensional or
three-dimensional CFD model. The simulator therefore provides two reduced-order
closures. This section describes the inexpensive integral model used by the
\emph{Quick screening} optimization mode. The nozzle core remains
quasi-one-dimensional, while the integrated mass-flow and momentum-flow
deficits produced by the unresolved boundary layer are transported along the
nozzle wall. Final evaluation uses the BLIMP-inspired profile marcher developed
in Section~\ref{sec:blimp_cebeci_smith}.

The general conservative transport equation can be written as

\begin{equation}
    \frac{\partial\boldsymbol{U}}{\partial t}
    +
    \nabla\cdot\boldsymbol{F}^{c}
    =
    \nabla\cdot\boldsymbol{F}^{d}
    +
    \boldsymbol{S},
    \label{eq:general_transport_BL}
\end{equation}

where \(\boldsymbol{U}\) is the vector of conserved variables,
\(\boldsymbol{F}^{c}\) is the convective-flux tensor,
\(\boldsymbol{F}^{d}\) is the diffusive-flux tensor, and
\(\boldsymbol{S}\) contains the remaining source terms.

For conservation of linear momentum, the transported conserved variable is

\begin{equation}
    \boldsymbol{U}_m=\rho\boldsymbol{u}.
\end{equation}

Neglecting body forces, the local momentum equation is

\begin{equation}
    \frac{\partial(\rho\boldsymbol{u})}{\partial t}
    +
    \nabla\cdot
    \left(
        \rho\boldsymbol{u}\otimes\boldsymbol{u}
    \right)
    =
    -\nabla p
    +
    \nabla\cdot\boldsymbol{\tau},
    \label{eq:momentum_transport_vector}
\end{equation}

where \(p\) is the static pressure and \(\boldsymbol{\tau}\) is the total
shear-stress tensor. For turbulent flow, \(\boldsymbol{\tau}\) represents the
combined molecular and modeled turbulent transport of momentum.

% ---------------------------------------------------------
\subsection{Coordinate System and Quasi-One-Dimensional Core}
% ---------------------------------------------------------

The boundary layer is described using the wall-fitted coordinates \(s\) and
\(n\). The coordinate \(s\) is measured tangentially along the nozzle wall,
whereas \(n\) is measured locally normal to the wall and towards the nozzle
core.

The nozzle contour is originally defined in axial coordinates by

\begin{equation}
    r=r(x).
\end{equation}

The differential wall distance is

\begin{equation}
    ds
    =
    \sqrt{dx^2+dr^2}
    =
    \sqrt{
        1+
        \left(
            \frac{dr}{dx}
        \right)^2
    }\,dx.
    \label{eq:wall_arc_length}
\end{equation}

Consequently,

\begin{equation}
    s(x)
    =
    \int_{x_0}^{x}
    \sqrt{
        1+
        \left(
            \frac{dr}{d\xi}
        \right)^2
    }\,d\xi.
    \label{eq:wall_coordinate}
\end{equation}

On the discrete nozzle grid,

\begin{equation}
    s_i
    =
    s_{i-1}
    +
    \sqrt{
        (x_i-x_{i-1})^2+
        (r_i-r_{i-1})^2
    }.
    \label{eq:discrete_wall_coordinate}
\end{equation}

The local wall radius is therefore

\begin{equation}
    R_w(s)=r\bigl(x(s)\bigr),
    \label{eq:local_wall_radius}
\end{equation}

and the corresponding wetted perimeter is

\begin{equation}
    \boxed{
    P_w(s)
    =
    2\pi R_w(s)
    =
    2\pi r\bigl(x(s)\bigr)
    }.
    \label{eq:wetted_perimeter_BL}
\end{equation}

The quasi-one-dimensional inviscid model provides the core-flow properties

\begin{equation}
    M(x),\quad
    T(x),\quad
    p(x),\quad
    \rho(x),\quad
    \gamma(x).
\end{equation}

These properties are mapped onto the wall coordinate through the one-to-one
correspondence \(x_i\leftrightarrow s_i\):

\begin{equation}
    \left\{
        M(x_i),T(x_i),p(x_i),\rho(x_i),\gamma(x_i)
    \right\}
    \longrightarrow
    \left\{
        M_e(s_i),T_e(s_i),p_e(s_i),\rho_e(s_i),\gamma_e(s_i)
    \right\}.
    \label{eq:edge_property_mapping}
\end{equation}

The subscript \(e\) denotes a property evaluated at the outer edge of the
boundary layer. In particular, the edge velocity is approximated by the local
quasi-one-dimensional core velocity:

\begin{equation}
    \boxed{
    U_e(s_i)
    \approx
    U_{\mathrm{Q1D}}(x_i)
    =
    M(x_i)
    \sqrt{
        \gamma(x_i)R_g(x_i)T(x_i)
    }
    },
    \qquad
    x_i\leftrightarrow s_i,
    \label{eq:Q1D_edge_velocity}
\end{equation}

where \(R_g\) is the specific gas constant and must not be confused with the
wall radius \(R_w\).

The edge velocity \(U_e(s)\) is therefore a one-dimensional quantity. It is
not a velocity profile across the boundary layer.

Figure~\ref{fig:core_BL_coordinates} illustrates the relation between the
quasi-one-dimensional core and the local wall-fitted coordinate system.


\begin{figure}[htbp]
\centering
\begin{tikzpicture}[
    scale=1.0,
    >=Latex,
    every node/.style={font=\small}
]

% Centreline
\draw[dash pattern=on 5pt off 3pt,gray]
    (-0.3,0) -- (10.6,0)
    node[right,black] {\(x\)};

% Inviscid core
\fill[green!8]
    (0,0)
    -- plot[smooth] coordinates {
        (0,2.60)
        (1.5,2.40)
        (2.8,1.60)
        (3.8,0.80)
        (5.0,1.00)
        (6.5,1.40)
        (8.2,1.75)
        (10.0,1.95)
    }
    -- (10.0,0)
    -- cycle;

% Boundary-layer region
\fill[orange!23]
    plot[smooth] coordinates {
        (0,3.00)
        (1.5,2.80)
        (2.8,1.90)
        (3.8,1.00)
        (5.0,1.25)
        (6.5,1.70)
        (8.2,2.10)
        (10.0,2.35)
    }
    --
    plot[smooth] coordinates {
        (10.0,1.95)
        (8.2,1.75)
        (6.5,1.40)
        (5.0,1.00)
        (3.8,0.80)
        (2.8,1.60)
        (1.5,2.40)
        (0,2.60)
    }
    -- cycle;

% Physical wall
\draw[very thick,blue!70!black]
    plot[smooth] coordinates {
        (0,3.00)
        (1.5,2.80)
        (2.8,1.90)
        (3.8,1.00)
        (5.0,1.25)
        (6.5,1.70)
        (8.2,2.10)
        (10.0,2.35)
    };

% BL edge
\draw[very thick,dashed,orange!80!black]
    plot[smooth] coordinates {
        (0,2.60)
        (1.5,2.40)
        (2.8,1.60)
        (3.8,0.80)
        (5.0,1.00)
        (6.5,1.40)
        (8.2,1.75)
        (10.0,1.95)
    };

\node[blue!70!black] at (1.25,3.25)
    {physical wall \(R_w(s)\)};

\node[orange!80!black] at (11,2.08)
    {boundary-layer edge};

\node[align=center] at (2.1,0.6)
    {quasi-1D core\\
    \(U_e(s),\,M_e(s),\,T_e(s),\,p_e(s)\)};

% Selected wall point
\coordinate (W) at (6.5,1.70);
\coordinate (E) at (6.5,1.40);

% Local tangent
\draw[->,very thick] (W) -- ++(1.35,0.4)
    node[above] {\(s\)};

% Local normal
\draw[->,very thick] (W) -- ++(-0.35,-1.30)
    node[below] {\(n\)};

% Local radius
\draw[<->,thick]
    (6.5,0.04) -- (6.5,1.70)
    node[midway,right] {\(R_w(s)\)};

% Edge velocity
\draw[->,very thick,green!50!black]
    (6.5,1.36) -- ++(1.25,0.34)
    node[below] {\(U_e(s)\)};

\end{tikzpicture}
\caption{Quasi-one-dimensional core and wall-fitted coordinates used by the
integral boundary-layer model.}
\label{fig:core_BL_coordinates}
\end{figure}

No additional resolved spatial dimension is introduced by the mapping
\(x\leftrightarrow s\). The inviscid core remains quasi-one-dimensional, and
the transported integral boundary-layer variables are also functions of the
single coordinate \(s\).

% ---------------------------------------------------------
\subsection{Velocity Components Inside the Boundary Layer}
% ---------------------------------------------------------

The actual local velocity field inside the boundary layer contains tangential
and wall-normal components:

\begin{equation}
    \boxed{
    \boldsymbol{u}(s,n)
    =
    u_s(s,n)\boldsymbol{e}_s
    +
    u_n(s,n)\boldsymbol{e}_n
    }.
    \label{eq:BL_velocity_decomposition}
\end{equation}

Here,

\begin{itemize}
    \item \(u_s(s,n)\) is the velocity component tangent to the wall;
    \item \(u_n(s,n)\) is the velocity component normal to the wall;
    \item \(U_e(s)\) is the one-dimensional tangential velocity at the
    boundary-layer edge.
\end{itemize}

At a stationary, impermeable, no-slip wall,

\begin{equation}
    \boxed{
    u_s(s,0)=0,
    \qquad
    u_n(s,0)=0
    }.
    \label{eq:BL_wall_conditions}
\end{equation}

At the outer edge of the boundary layer,

\begin{equation}
    u_s\bigl(s,\delta(s)\bigr)\simeq U_e(s),
    \qquad
    \tau_{sn}\bigl(s,\delta(s)\bigr)\simeq0.
    \label{eq:BL_edge_conditions}
\end{equation}

The wall-normal component is zero at the impermeable wall but is generally
nonzero inside the boundary layer:

\begin{equation}
    u_n(s,n)\neq0,
    \qquad
    0<n\leq\delta(s).
    \label{eq:normal_velocity_nonzero}
\end{equation}

Its magnitude is normally much smaller than the tangential velocity:

\begin{equation}
    |u_n|\ll|u_s|.
    \label{eq:BL_velocity_scaling}
\end{equation}

Nevertheless, the normal convection term cannot be discarded because the
normal velocity gradient is large inside a thin boundary layer. In general,

\begin{equation}
    u_n\frac{\partial u_s}{\partial n}
\end{equation}

can be of the same asymptotic order as

\begin{equation}
    u_s\frac{\partial u_s}{\partial s}.
\end{equation}

The present model does not explicitly resolve
\(u_s(s,n)\), \(u_n(s,n)\), \(T(s,n)\), or \(\rho(s,n)\). These
two-dimensional fields appear in the theoretical definitions and derivation,
but their integrated effects are represented by one-dimensional
boundary-layer variables.

% ---------------------------------------------------------
\subsection{Physical and Integral Boundary-Layer Thicknesses}
% ---------------------------------------------------------

Three different boundary-layer thicknesses are employed:

\begin{equation}
    \boxed{
    \delta(s),\qquad
    \delta^*(s),\qquad
    \theta(s)
    }.
\end{equation}

Although all three quantities have dimensions of length, they represent
different physical effects.

\begin{figure}[htbp]
\centering
\begin{tikzpicture}[
    x=1.0cm,
    y=1.0cm,
    >=Latex,
    every node/.style={font=\small}
]

% ---------------------------------------------------------
% Coordinates
% ---------------------------------------------------------

\draw[->,thick] (-0.2,0) -- (14.7,0)
    node[below] {\(s\)};

\draw[->,thick] (0,-0.2) -- (0,4.65)
    node[left] {\(n\)};

% Wall
\draw[very thick] (0,0) -- (10.2,0)
    node[below left] {wall};

% ---------------------------------------------------------
% Boundary-layer region
% ---------------------------------------------------------

\fill[orange!18]
    (0,0)
    .. controls (2.8,1.10) and (7.0,2.55) ..
    (10.6,3.15)
    -- (10.6,0)
    -- cycle;

\draw[very thick,orange!80!black]
    (0,0)
    .. controls (2.8,1.10) and (7.0,2.55) ..
    (10.6,3.15);

\node[orange!80!black,rotate=13]
    at (6.85,2.53)
    {\(n=\delta(s)\)};

\node[orange!80!black]
    at (4.0,2.30)
    {boundary layer};

% ---------------------------------------------------------
% Velocity profile at s_1
% ---------------------------------------------------------

\draw[very thick,blue!70!black]
    (2.0,0)
    .. controls (2.20,0.15) and (2.95,0.70) ..
    (3.20,1.15);

\foreach \yy/\xx in {
    0.20/2.16,
    0.43/2.39,
    0.68/2.68,
    0.91/2.96,
    1.10/3.17
}{
    \draw[->,blue!70!black]
        (2.0,\yy) -- (\xx,\yy);
}

\draw[densely dashed,gray]
    (3.20,0) -- (3.20,1.15);

\draw[<->,thick]
    (1.65,0) -- (1.65,1.15)
    node[midway,left] {\(\delta(s_1)\)};

\node[below] at (2.0,0) {\(s_1\)};

% ---------------------------------------------------------
% Velocity profile at s_2
% ---------------------------------------------------------

\draw[very thick,blue!70!black]
    (5.0,0)
    .. controls (5.25,0.20) and (6.12,1.13) ..
    (6.35,2.08);

\foreach \yy/\xx in {
    0.28/5.19,
    0.62/5.45,
    0.99/5.76,
    1.36/6.02,
    1.72/6.24,
    2.02/6.34
}{
    \draw[->,blue!70!black]
        (5.0,\yy) -- (\xx,\yy);
}

\draw[densely dashed,gray]
    (6.35,0) -- (6.35,2.08);

\draw[<->,thick]
    (4.62,0) -- (4.62,2.08)
    node[midway,left] {\(\delta(s_2)\)};

\node[below] at (5.0,0) {\(s_2\)};

% ---------------------------------------------------------
% Velocity profile at s_3
% ---------------------------------------------------------

\draw[very thick,blue!70!black]
    (8.0,0)
    .. controls (8.30,0.22) and (9.28,1.64) ..
    (9.50,2.88);

\foreach \yy/\xx in {
    0.32/8.20,
    0.73/8.49,
    1.17/8.82,
    1.61/9.08,
    2.04/9.30,
    2.43/9.43,
    2.80/9.50
}{
    \draw[->,blue!70!black]
        (8.0,\yy) -- (\xx,\yy);
}

\draw[densely dashed,gray]
    (9.50,0) -- (9.50,2.88);

\node[below] at (8.0,0) {\(s_3\)};

% Edge velocity
\draw[->,very thick,blue!70!black]
    (8.0,3.54) -- (9.50,3.54)
    node[right] {\(U_e(s_3)\)};

% ---------------------------------------------------------
% Physical and equivalent thicknesses at s_3
% Values are schematic, not quantitatively scaled
% ---------------------------------------------------------

% Momentum thickness theta
% Displacement thickness delta*  (fica primeiro / mais à esquerda)
\draw[densely dashed,red!70!black]
    (8.0,0.78) -- (11.10,0.78);

\draw[<->,very thick,red!70!black]
    (11.10,0) -- (11.10,0.78);

\node[
    right,
    red!70!black,
    align=left
] at (11.17,1.12)
    {\(\delta^*(s_3)\)\\[-1mm]
     \scriptsize displacement\\ \scriptsize thickness};


% Momentum thickness theta  (fica depois / mais à direita)
\draw[densely dashed,green!50!black]
    (8.0,0.52) -- (12.85,0.52);

\draw[<->,very thick,green!50!black]
    (12.85,0) -- (12.85,0.52);

\node[
    right,
    green!50!black,
    align=left
] at (12.92,0.46)
    {\(\theta(s_3)\)\\[-1mm]
     \scriptsize momentum thickness};

% Physical thickness delta
\draw[densely dashed,orange!80!black]
    (8.0,2.88) -- (10.05,2.88);

\draw[<->,very thick,orange!80!black]
    (10.05,0) -- (10.05,2.88);

\node[
    right,
    orange!80!black,
    align=left
] at (10.12,2.40)
    {\(\delta(s_3)\)\\[-1mm]
     \scriptsize physical thickness};

% Ordering
\node[
    draw,
    rounded corners,
    fill=white,
    align=center,
    inner sep=4pt
] at (10.55,4.25)
    {\(\theta(s_3)<\delta^*(s_3)<\delta(s_3)\)};

% Clarification

% ---------------------------------------------------------
% No-slip condition
% ---------------------------------------------------------

\draw[->,thick]
    (1.15,-0.60) -- (1.96,0.04);

\node[align=center,left]
    at (1.15,-0.70)
    {no slip\\\(u_s(s,0)=0\)};

% Profile label
\node[
    blue!70!black,
    align=center
] at (6.8,0.45)
    {\scriptsize tangential velocity\\ \scriptsize profiles \(u_s(s,n)\)};

\end{tikzpicture}

\caption{Development of the tangential velocity profile along the nozzle
wall and comparison of the conventional boundary-layer thickness
\(\delta\), displacement thickness \(\delta^*\), and momentum thickness
\(\theta\) at station \(s_3\). The coordinate \(n\) is normal to the wall,
whereas \(s\) follows the wall in the streamwise direction. The displayed
relative magnitudes of \(\theta\) and \(\delta^*\) are schematic and are not
drawn to scale.}

\label{fig:BL_velocity_growth}
\end{figure}
\subsubsection{Representative physical thickness}

The velocity approaches \(U_e\) asymptotically, and therefore the boundary
layer does not possess an exact physical edge. A common engineering convention
is to identify its thickness with the wall-normal location at which

\begin{equation}
    \frac{u_s}{U_e}\simeq0.99.
    \label{eq:delta_99_definition}
\end{equation}

The quantity \(\delta(s)\) is consequently interpreted as a representative
physical thickness of the velocity-gradient region. In the profile closure
introduced below, it is treated as the outer scale of the assumed velocity
profile and is approximately identified with the conventional
\(\delta_{99}\).

\subsubsection{Compressible displacement thickness}

Inside the boundary layer, the local mass flux \(\rho u_s\) is smaller than
the mass flux \(\rho_eU_e\) that would exist if the inviscid edge state
extended uniformly to the physical wall.

The mass-flow deficit per unit wetted perimeter is

\begin{equation}
    \dot{m}'_{\mathrm{def}}(s)
    =
    \int_0^{\delta(s)}
    \left[
        \rho_e(s)U_e(s)
        -
        \rho(s,n)u_s(s,n)
    \right]dn.
    \label{eq:mass_flow_deficit}
\end{equation}

Factoring out the edge mass flux gives

\begin{equation}
    \dot{m}'_{\mathrm{def}}(s)
    =
    \rho_e(s)U_e(s)
    \int_0^{\delta(s)}
    \left[
        1-
        \frac{
            \rho(s,n)u_s(s,n)
        }{
            \rho_e(s)U_e(s)
        }
    \right]dn.
\end{equation}

The compressible displacement thickness is therefore defined as

\begin{equation}
    \boxed{
    \delta^*(s)
    =
    \int_0^{\delta(s)}
    \left[
        1-
        \frac{
            \rho(s,n)u_s(s,n)
        }{
            \rho_e(s)U_e(s)
        }
    \right]dn
    }.
    \label{eq:compressible_displacement_thickness}
\end{equation}

It follows that

\begin{equation}
    \boxed{
    \dot{m}'_{\mathrm{def}}(s)
    =
    \rho_e(s)U_e(s)\delta^*(s)
    }.
    \label{eq:mass_deficit_displacement}
\end{equation}

The displacement thickness is the thickness of a hypothetical region carrying
zero mass flow that would produce the same mass-flow deficit as the actual
boundary layer.

Equivalently, the inviscid core behaves approximately as if the solid wall
were displaced towards the nozzle axis by \(\delta^*\). Under the thin-layer
approximation, the effective core radius is

\begin{equation}
    \boxed{
    R_{\mathrm{eff}}(s)
    =
    R_w(s)-\delta^*(s)
    }.
    \label{eq:effective_radius_displacement}
\end{equation}

The displacement thickness is not the physical location of the
boundary-layer edge. Generally,

\begin{equation}
    \delta^*(s)<\delta(s).
    \label{eq:displacement_less_than_delta}
\end{equation}

\subsubsection{Compressible momentum thickness}

The reduced tangential velocity also produces a deficit in streamwise
momentum flux. The momentum-flux deficit per unit wetted perimeter is

\begin{equation}
    D_m'(s)
    =
    \int_0^{\delta(s)}
    \rho(s,n)u_s(s,n)
    \left[
        U_e(s)-u_s(s,n)
    \right]dn.
    \label{eq:momentum_deficit_integral}
\end{equation}

Introducing the edge quantities gives

\begin{equation}
\begin{split}
    D_m'(s)
    &=
    \rho_e(s)U_e^2(s)
    \int_0^{\delta(s)}
    \frac{
        \rho(s,n)u_s(s,n)
    }{
        \rho_e(s)U_e(s)
    }
    \left[
        1-
        \frac{
            u_s(s,n)
        }{
            U_e(s)
        }
    \right]dn.
\end{split}
\label{eq:momentum_deficit_normalized}
\end{equation}

The compressible momentum thickness is defined as

\begin{equation}
    \boxed{
    \theta(s)
    =
    \int_0^{\delta(s)}
    \frac{
        \rho(s,n)u_s(s,n)
    }{
        \rho_e(s)U_e(s)
    }
    \left[
        1-
        \frac{
            u_s(s,n)
        }{
            U_e(s)
        }
    \right]dn
    }.
    \label{eq:compressible_momentum_thickness}
\end{equation}

Consequently,

\begin{equation}
    \boxed{
    D_m'(s)
    =
    \rho_e(s)U_e^2(s)\theta(s)
    },
    \qquad
    [D_m']=\mathrm{N\,m^{-1}}.
    \label{eq:momentum_deficit_per_perimeter}
\end{equation}

The momentum thickness is the thickness of a hypothetical region producing
the same momentum-flux deficit as the actual velocity and density profiles. It
is not itself momentum and does not represent a physical interface inside the
boundary layer.

For the complete circumference of an axisymmetric nozzle,

\begin{equation}
    \boxed{
    D_m(s)
    =
    P_w(s)\rho_e(s)U_e^2(s)\theta(s)
    }
    \label{eq:axisymmetric_momentum_deficit}
\end{equation}

or, using \(P_w=2\pi R_w\),

\begin{equation}
    \boxed{
    D_m(s)
    =
    2\pi R_w(s)\rho_e(s)U_e^2(s)\theta(s)
    },
    \qquad
    [D_m]=\mathrm{N}.
    \label{eq:axisymmetric_momentum_deficit_radius}
\end{equation}

The three thicknesses can therefore be summarized as

\begin{align}
    \delta
    &:
    &&\text{representative physical extent of the velocity-gradient region},
    \\
    \delta^*
    &:
    &&\text{equivalent thickness associated with mass-flow deficit},
    \\
    \theta
    &:
    &&\text{equivalent thickness associated with momentum-flux deficit}.
\end{align}

For a conventional attached turbulent boundary layer, the typical ordering is

\begin{equation}
    \boxed{
    \theta(s)<\delta^*(s)<\delta(s)
    }.
    \label{eq:BL_thickness_ordering}
\end{equation}

The ratio

\begin{equation}
    \boxed{
    H(s)
    =
    \frac{\delta^*(s)}{\theta(s)}
    }
    \label{eq:shape_factor}
\end{equation}

is the compressible boundary-layer shape factor. It provides an integral
measure of the shapes of the velocity and density profiles.

Figure~\ref{fig:BL_thicknesses} illustrates the different meanings of
\(\delta\), \(\delta^*\), and \(\theta\).

\begin{figure}[htbp]
\centering
\begin{tikzpicture}[
    x=1.05cm,
    y=1.0cm,
    >=Latex,
    every node/.style={font=\small}
]

% Axes
\draw[->,thick] (0,0) -- (4.6,0)
    node[right] {\(u_s/U_e\)};
\draw[->,thick] (0,0) -- (0,4.25)
    node[above] {\(n\)};

% Edge velocity
\draw[dashed,gray]
    (3.7,0) -- (3.7,3.55);
\node[below] at (3.7,0) {\(1\)};

% Velocity profile
\draw[very thick,blue!70!black]
    (0,0)
    .. controls (0.8,0.18) and (2.8,1.35) ..
    (3.42,2.70)
    .. controls (3.58,3.05) and (3.66,3.35) ..
    (3.70,3.55);

% Deficit area
\fill[orange!25,opacity=0.65]
    (3.7,0)
    -- (0,0)
    .. controls (0.8,0.18) and (2.8,1.35) ..
    (3.42,2.70)
    .. controls (3.58,3.05) and (3.66,3.35) ..
    (3.70,3.55)
    -- cycle;

% Delta
\draw[<->,thick]
    (-0.43,0) -- (-0.43,3.55)
    node[midway,left] {\(\delta\)};
\draw[dashed] (0,3.55) -- (3.70,3.55);

\node[align=center,orange!80!black]
    at (2.8,0.65)
    {mass and\\ momentum\\flux deficits};

% Equivalent bars
\draw[->,thick] (6.0,0) -- (11.8,0);
\draw[->,thick] (6.0,0) -- (6.0,4.25)
    node[above] {equivalent length};

\fill[blue!15] (6.8,0) rectangle (7.8,3.55);
\draw[thick,blue!70!black] (6.8,0) rectangle (7.8,3.55);
\node[below] at (7.3,0) {\(\delta\)};
\node[align=center] at (7.3,3.95) {physical\\extent};

\fill[orange!30] (8.7,0) rectangle (9.7,1.70);
\draw[thick,orange!80!black] (8.7,0) rectangle (9.7,1.70);
\node[below] at (9.2,0) {\(\delta^*\)};
\node[align=center] at (9.2,2.08) {mass-flow\\deficit};

\fill[green!30] (10.6,0) rectangle (11.6,1.10);
\draw[thick,green!50!black] (10.6,0) rectangle (11.6,1.10);
\node[below] at (11.1,0) {\(\theta\)};
\node[align=center] at (11.1,1.48) {momentum-flux\\deficit};

\end{tikzpicture}

\caption{Physical interpretation of the representative boundary-layer
thickness \(\delta\), displacement thickness \(\delta^*\), and momentum
thickness \(\theta\). The relative dimensions are schematic and not to
scale.}
\label{fig:BL_thicknesses}
\end{figure}

% ---------------------------------------------------------
\subsection{Normalized Profile Closure}
% ---------------------------------------------------------

The definitions of \(\delta^*\) and \(\theta\) depend on the wall-normal
profiles

\begin{equation}
    u_s(s,n),\qquad
    T(s,n),\qquad
    \rho(s,n),
\end{equation}

and on the physical thickness \(\delta(s)\). These quantities are not provided
by the quasi-one-dimensional core model.

To evaluate the shape factor without resolving the complete two-dimensional
boundary layer, normalized velocity, temperature, and density profiles are
introduced.

\subsubsection{Normalized wall-normal coordinate}

The dimensionless coordinate

\begin{equation}
    \boxed{
    \eta
    =
    \frac{n}{\delta(s)}
    },
    \qquad
    0\leq\eta\leq1,
    \label{eq:normalized_wall_coordinate}
\end{equation}

maps the boundary layer at every streamwise station onto the same normalized
interval.

Thus,

\begin{align}
    \eta=0
    &\quad\Longleftrightarrow\quad
    n=0,
    &&\text{wall},
    \\
    \eta=1
    &\quad\Longleftrightarrow\quad
    n=\delta(s),
    &&\text{representative boundary-layer edge}.
\end{align}

The dimensional and normalized coordinates are related by

\begin{equation}
    n=\eta\delta(s),
    \qquad
    dn=\delta(s)\,d\eta.
    \label{eq:normal_coordinate_transformation}
\end{equation}

The coordinate \(\eta\) is not an additional resolved spatial dimension. It is
an auxiliary quadrature coordinate used to describe the assumed shapes of the
unresolved profiles.

The normalization is useful because \(\delta(s)\) is not initially known.
Although the dimensional values of \(\delta^*\) and \(\theta\) cannot yet be
calculated from their definitions, their ratios to \(\delta\) can be
calculated using normalized profiles.

\subsubsection{One-seventh-power velocity profile}

The tangential velocity profile is approximated using the empirical turbulent
one-seventh-power law:

\begin{equation}
    \boxed{
    f(\eta)
    \equiv
    \frac{u_s(s,n)}{U_e(s)}
    =
    \eta^{1/7}
    }.
    \label{eq:one_seventh_velocity_profile}
\end{equation}

The profile satisfies the idealized boundary conditions

\begin{equation}
    f(0)=0,
    \qquad
    f(1)=1.
\end{equation}

The first condition represents no slip at the wall. The second condition
matches the assumed profile to the edge velocity. The conventional
\(99\%\) definition of \(\delta\) and the condition \(f(1)=1\) are treated as
equivalent engineering representations of the boundary-layer edge.

The one-seventh-power law prescribes only the normalized shape of the velocity
profile. It does not determine the dimensional thickness \(\delta(s)\).

Figure~\ref{fig:normalized_BL_profiles} illustrates the normalized coordinate
and the assumed velocity and temperature profiles.

\begin{figure}[htbp]
\centering
\begin{tikzpicture}[
    x=1.1cm,
    y=1.0cm,
    >=Latex,
    every node/.style={font=\small}
]

% Coordinate axes
\draw[->,thick] (0,0) -- (4.4,0)
    node[below] {\(u_s/U_e,\;T/T_{\mathrm{rec}}\)};
\draw[->,thick] (0,0) -- (0,4.3)
    node[above] {\(\eta=n/\delta\)};

% Eta boundaries
\draw[dashed,gray] (0,3.65) -- (4.0,3.65);
\node[left] at (0,3.65) {\(1\)};
\node[left] at (0,0) {\(0\)};

% Velocity profile
\draw[very thick,blue!70!black]
    (0,0)
    .. controls (1.20,0.10) and (3.25,1.70) ..
    (3.75,3.65);

\node[blue!70!black,align=center]
    at (2.65,0.35)
    {\(\displaystyle
      \frac{u_s}{U_e}=\eta^{1/7}\)};

% Temperature profile
\draw[very thick,red!70!black]
    (3.75,0)
    .. controls (3.65,1.20) and (3.25,2.60) ..
    (2.65,3.65);

\node[red!70!black,align=center]
    at (1.75,2.85)
    {\(T(s,\eta)\)\\from Walz};

% Wall and edge
\node[align=left] at (7.4,0.25)
    {\(\eta=0:\quad u_s=0,\quad T=T_{\mathrm{rec}}\)};

\node[align=left] at (7.4,3.60)
    {\(\eta=1:\quad u_s=U_e,\quad T=T_e\)};

\draw[->,thick] (4.0,0.25) -- (5.0,0.25);
\draw[->,thick] (4.0,3.60) -- (5.0,3.60);

\end{tikzpicture}

\caption{Normalized wall-normal coordinate and assumed profiles for an
adiabatic compressible turbulent boundary layer. The diagram is schematic.}
\label{fig:normalized_BL_profiles}
\end{figure}

\subsubsection{Recovery temperature and adiabatic wall}

The static edge temperature \(T_e(s)\) does not represent the complete
temperature distribution across the boundary layer. In compressible flow,
part of the kinetic energy is recovered as internal energy near the wall.

The corresponding recovery temperature is

\begin{equation}
    \boxed{
    T_{\mathrm{rec}}(s)
    =
    T_e(s)
    \left[
        1+
        r_{\mathrm{rec}}(s)
        \frac{\gamma_e(s)-1}{2}
        M_e^2(s)
    \right]
    },
    \label{eq:recovery_temperature}
\end{equation}

where \(r_{\mathrm{rec}}\) is the recovery factor. For a turbulent boundary
layer,

\begin{equation}
    \boxed{
    r_{\mathrm{rec}}(s)
    \simeq
    Pr_e^{1/3}(s)
    }.
    \label{eq:turbulent_recovery_factor}
\end{equation}

The recovery temperature is the temperature approached by an adiabatic wall
under the local edge-flow conditions. It is generally bounded by

\begin{equation}
    T_e(s)
    \leq
    T_{\mathrm{rec}}(s)
    \leq
    T_0(s),
\end{equation}

where \(T_0\) is the local stagnation temperature.

The present model assumes an adiabatic wall:

\begin{equation}
    q_w''=0.
\end{equation}

Consequently,

\begin{equation}
    \boxed{
    T_w(s)=T_{\mathrm{rec}}(s)
    }.
    \label{eq:quick_adiabatic_wall_temperature}
\end{equation}

The wall temperature is therefore not an independent prescribed input in the
adiabatic boundary-layer closure.

\subsubsection{Walz temperature--velocity relation}

The density profile required by \(\delta^*\) and \(\theta\) depends on the
temperature across the boundary layer. Since the present model does not solve
a wall-normal energy equation, the temperature profile is approximated using
the Walz temperature--velocity relation.

For a general wall temperature, the Walz relation is

\begin{equation}
    \boxed{
    T(s,\eta)
    =
    T_w(s)
    +
    \left[
        T_{\mathrm{rec}}(s)-T_w(s)
    \right]f(\eta)
    -
    \left[
        T_{\mathrm{rec}}(s)-T_e(s)
    \right]f^2(\eta)
    }.
    \label{eq:general_walz_relation}
\end{equation}

There is no single quantity referred to here as the ``Walz temperature.''
Equation~\eqref{eq:general_walz_relation} provides the complete approximate
temperature profile

\begin{equation}
    T=T(s,\eta)
\end{equation}

as a function of the local normalized velocity \(f=u_s/U_e\).

Under the adiabatic-wall condition \(T_w=T_{\mathrm{rec}}\), the linear term
vanishes and the relation reduces to

\begin{equation}
    \boxed{
    T(s,\eta)
    =
    T_{\mathrm{rec}}(s)
    -
    \left[
        T_{\mathrm{rec}}(s)-T_e(s)
    \right]f^2(\eta)
    }.
    \label{eq:adiabatic_walz_relation}
\end{equation}

Using \(f(\eta)=\eta^{1/7}\),

\begin{equation}
    \boxed{
    T(s,\eta)
    =
    T_{\mathrm{rec}}(s)
    -
    \left[
        T_{\mathrm{rec}}(s)-T_e(s)
    \right]\eta^{2/7}
    }.
    \label{eq:adiabatic_walz_eta}
\end{equation}

The limiting values are

\begin{align}
    T(s,0)
    &=
    T_{\mathrm{rec}}(s)
    =
    T_w(s),
    \\
    T(s,1)
    &=
    T_e(s).
\end{align}

The relation therefore describes the temperature variation from the recovery
temperature at the adiabatic wall to the static edge temperature at the outer
boundary of the layer.

\subsubsection{Normalized density profile}

The thin-boundary-layer approximation implies that the static pressure is
approximately uniform in the wall-normal direction:

\begin{equation}
    \boxed{
    p(s,n)\simeq p_e(s)
    },
    \qquad
    \frac{\partial p}{\partial n}\simeq0.
    \label{eq:normal_pressure_uniformity}
\end{equation}

This does not imply that the density is uniform across the boundary layer.
Since

\begin{equation}
    T=T(s,n),
\end{equation}

the density also depends on the wall-normal coordinate:

\begin{equation}
    \rho=\rho(s,n).
\end{equation}

Assuming ideal-gas behavior and an approximately uniform specific gas
constant across the thin layer,

\begin{align}
    \rho(s,\eta)
    &=
    \frac{p_e(s)}
         {R_g(s)T(s,\eta)},
    \\
    \rho_e(s)
    &=
    \frac{p_e(s)}
         {R_g(s)T_e(s)}.
\end{align}

Dividing these expressions gives the normalized density profile

\begin{equation}
    \boxed{
    \frac{\rho(s,\eta)}{\rho_e(s)}
    =
    \frac{T_e(s)}{T(s,\eta)}
    }.
    \label{eq:normalized_density_profile}
\end{equation}

The assumed profile closure can therefore be summarized as

\begin{equation}
    \boxed{
    \eta
    \longrightarrow
    \frac{u_s}{U_e}
    \longrightarrow
    T(s,\eta)
    \longrightarrow
    \frac{\rho(s,\eta)}{\rho_e(s)}
    }.
    \label{eq:profile_closure_sequence}
\end{equation}

The one-seventh-power law provides the normalized velocity profile, the Walz
relation provides the corresponding temperature profile, and the ideal-gas
relation provides the density profile required by the compressible integral
thicknesses.

\subsubsection{Normalized integral thicknesses}

Using

\begin{equation}
    n=\eta\delta(s),
    \qquad
    dn=\delta(s)\,d\eta,
\end{equation}

the displacement thickness becomes

\begin{equation}
\begin{split}
    \delta^*(s)
    &=
    \delta(s)
    \int_0^1
    \left[
        1-
        \frac{\rho(s,\eta)}{\rho_e(s)}
        f(\eta)
    \right]d\eta.
\end{split}
\label{eq:normalized_displacement_thickness}
\end{equation}

The dimensionless displacement-thickness ratio is therefore

\begin{equation}
    \boxed{
    I_{\delta^*}(s)
    \equiv
    \frac{\delta^*(s)}{\delta(s)}
    =
    \int_0^1
    \left[
        1-
        \frac{\rho(s,\eta)}{\rho_e(s)}
        f(\eta)
    \right]d\eta
    }.
    \label{eq:displacement_thickness_ratio}
\end{equation}

Similarly, the momentum thickness becomes

\begin{equation}
\begin{split}
    \theta(s)
    &=
    \delta(s)
    \int_0^1
    \frac{\rho(s,\eta)}{\rho_e(s)}
    f(\eta)
    \left[
        1-f(\eta)
    \right]d\eta.
\end{split}
\label{eq:normalized_momentum_thickness}
\end{equation}

The corresponding dimensionless ratio is

\begin{equation}
    \boxed{
    I_\theta(s)
    \equiv
    \frac{\theta(s)}{\delta(s)}
    =
    \int_0^1
    \frac{\rho(s,\eta)}{\rho_e(s)}
    f(\eta)
    \left[
        1-f(\eta)
    \right]d\eta
    }.
    \label{eq:momentum_thickness_ratio}
\end{equation}

At this stage, the dimensional boundary-layer thickness \(\delta(s)\) remains
unknown. Consequently, these profile integrations do not directly provide
\(\delta^*\) or \(\theta\) in metres.

However, the unknown thickness cancels in their ratio:

\begin{equation}
\begin{split}
    H(s)
    &=
    \frac{\delta^*(s)}{\theta(s)}
    =
    \frac{
        \delta(s)I_{\delta^*}(s)
    }{
        \delta(s)I_\theta(s)
    }.
\end{split}
\end{equation}

Therefore,

\begin{equation}
    \boxed{
    H(s)
    =
    \frac{I_{\delta^*}(s)}
         {I_\theta(s)}
    }
    \label{eq:shape_factor_ratio}
\end{equation}

or, explicitly,

\begin{equation}
    \boxed{
    H(s)
    =
    \frac{
        \displaystyle
        \int_0^1
        \left[
            1-
            \frac{\rho}{\rho_e}f
        \right]d\eta
    }{
        \displaystyle
        \int_0^1
        \frac{\rho}{\rho_e}
        f(1-f)\,d\eta
    }
    }.
    \label{eq:shape_factor_profile_closure}
\end{equation}

This cancellation is the principal reason for introducing the normalized
coordinate \(\eta\). The shape factor can be calculated before the dimensional
boundary-layer thickness is known.

Numerically, the interval \(0\leq\eta\leq1\) is divided into a set of
quadrature points. At every streamwise station \(s_i\), the model evaluates

\begin{equation}
    f(\eta_j),\qquad
    T(s_i,\eta_j),\qquad
    \frac{\rho(s_i,\eta_j)}{\rho_e(s_i)},
\end{equation}

and integrates Eqs.~\eqref{eq:displacement_thickness_ratio} and
\eqref{eq:momentum_thickness_ratio} using the trapezoidal rule.

Once \(\theta(s)\) has been obtained from the integral momentum equation, the
remaining thicknesses can be recovered from

\begin{equation}
    \boxed{
    \delta^*(s)=H(s)\theta(s)
    }
    \label{eq:delta_star_from_H_theta}
\end{equation}

and

\begin{equation}
    \boxed{
    \delta(s)
    =
    \frac{\theta(s)}
         {I_\theta(s)}
    }.
    \label{eq:physical_delta_from_theta}
\end{equation}

% ---------------------------------------------------------
\subsection{Local Conservative Momentum Equation}
% ---------------------------------------------------------

For statistically steady flow,

\begin{equation}
    \frac{\partial(\rho\boldsymbol{u})}{\partial t}=0.
    \label{eq:steady_momentum_condition}
\end{equation}

The conservative momentum equation therefore reduces to

\begin{equation}
    \nabla\cdot
    \left(
        \rho\boldsymbol{u}\otimes\boldsymbol{u}
    \right)
    =
    -\nabla p
    +
    \nabla\cdot\boldsymbol{\tau}.
    \label{eq:steady_conservative_momentum}
\end{equation}

Resolving this equation in the direction tangent to the wall gives

\begin{equation}
    \frac{\partial(\rho u_s^2)}{\partial s}
    +
    \frac{\partial(\rho u_nu_s)}{\partial n}
    =
    -\frac{\partial p}{\partial s}
    +
    \frac{\partial\tau_{ss}}{\partial s}
    +
    \frac{\partial\tau_{sn}}{\partial n}.
    \label{eq:local_tangential_momentum}
\end{equation}

The first two terms represent the tangential and normal convective transport
of tangential momentum. The last two terms represent the tangential and normal
diffusive transport of tangential momentum.

For a thin boundary layer,

\begin{equation}
    \delta\ll L,
    \qquad
    \frac{\partial}{\partial n}
    \gg
    \frac{\partial}{\partial s}.
    \label{eq:thin_BL_scale}
\end{equation}

Consequently, the streamwise diffusion of tangential momentum is much smaller
than its wall-normal diffusion:

\begin{equation}
    \frac{\partial\tau_{ss}}{\partial s}
    \ll
    \frac{\partial\tau_{sn}}{\partial n}.
    \label{eq:negligible_streamwise_diffusion}
\end{equation}

Using also \(p(s,n)\simeq p_e(s)\), the local tangential momentum equation
reduces to

\begin{equation}
    \boxed{
    \frac{\partial(\rho u_s^2)}{\partial s}
    +
    \frac{\partial(\rho u_nu_s)}{\partial n}
    =
    -\frac{dp_e}{ds}
    +
    \frac{\partial\tau_{sn}}{\partial n}
    }.
    \label{eq:reduced_tangential_momentum}
\end{equation}

The corresponding steady continuity equation is

\begin{equation}
    \nabla\cdot(\rho\boldsymbol{u})=0.
    \label{eq:steady_continuity_BL}
\end{equation}

For a thin axisymmetric boundary layer, it can be expressed approximately in
perimeter-weighted form as

\begin{equation}
    \boxed{
    \frac{\partial}{\partial s}
    \left[
        P_w(s)\rho u_s
    \right]
    +
    \frac{\partial}{\partial n}
    \left[
        P_w(s)\rho u_n
    \right]
    =0
    }.
    \label{eq:axisymmetric_BL_continuity}
\end{equation}

The normal velocity \(u_n\) is required by local mass conservation. It is zero
at the impermeable wall but is generally nonzero inside the boundary layer.
The integral model does not solve \(u_n(s,n)\) explicitly. Instead, continuity
is combined with the tangential momentum equation before integration in the
normal direction.

% ---------------------------------------------------------
\subsection{Integral Momentum Equation}
% ---------------------------------------------------------

The inviscid edge flow satisfies the streamwise Euler equation

\begin{equation}
    \boxed{
    \rho_e(s)U_e(s)\frac{dU_e}{ds}
    =
    -\frac{dp_e}{ds}
    }.
    \label{eq:edge_euler_equation}
\end{equation}

Equations~\eqref{eq:reduced_tangential_momentum},
\eqref{eq:axisymmetric_BL_continuity}, and
\eqref{eq:edge_euler_equation} are combined and integrated from the wall,
\(n=0\), to the boundary-layer edge, \(n=\delta(s)\).

During this integration, the normal convective transport containing \(u_n\)
reduces to boundary contributions. At the impermeable wall,

\begin{equation}
    u_n(s,0)=0,
\end{equation}

whereas the momentum-deficit integrand vanishes at the boundary-layer edge
because

\begin{equation}
    u_s\bigl(s,\delta(s)\bigr)\simeq U_e(s).
\end{equation}

The wall-normal diffusion term becomes

\begin{equation}
    \int_0^{\delta(s)}
    \frac{\partial\tau_{sn}}{\partial n}\,dn
    =
    \tau_{sn}\bigl(s,\delta(s)\bigr)
    -
    \tau_{sn}(s,0).
    \label{eq:integrated_normal_diffusion}
\end{equation}

Since the shear stress is negligible at the boundary-layer edge,

\begin{equation}
    \tau_{sn}\bigl(s,\delta(s)\bigr)\simeq0,
\end{equation}

the integrated normal diffusion is represented by the wall shear stress.
Defining its positive magnitude as

\begin{equation}
    \tau_w(s)
    =
    \left|
        \tau_{sn}(s,0)
    \right|,
    \label{eq:wall_shear_definition}
\end{equation}

the axisymmetric compressible integral momentum equation becomes

\begin{equation}
    \boxed{
    \frac{d}{ds}
    \left[
        P_w(s)\rho_e(s)U_e^2(s)\theta(s)
    \right]
    +
    P_w(s)\rho_e(s)U_e(s)
    \frac{dU_e}{ds}
    \delta^*(s)
    =
    P_w(s)\tau_w(s)
    }.
    \label{eq:axisymmetric_integral_momentum}
\end{equation}

Since

\begin{equation}
    D_m(s)
    =
    P_w(s)\rho_e(s)U_e^2(s)\theta(s),
\end{equation}

the integral equation can also be written as

\begin{equation}
    \boxed{
    \frac{dD_m}{ds}
    =
    P_w\tau_w
    -
    P_w\rho_eU_e
    \frac{dU_e}{ds}\delta^*
    }.
    \label{eq:momentum_deficit_transport}
\end{equation}

This equation shows that the transported integral quantity is the total
momentum-flux deficit \(D_m\), rather than \(\theta\) alone. The momentum
thickness is the equivalent length used to express this deficit through

\begin{equation}
    D_m=P_w\rho_eU_e^2\theta.
\end{equation}

Defining the skin-friction coefficient as

\begin{equation}
    \boxed{
    C_f(s)
    =
    \frac{2\tau_w(s)}
         {\rho_e(s)U_e^2(s)}
    }
    \label{eq:skin_friction_closure_problem}
\end{equation}

and using

\begin{equation}
    H(s)=\frac{\delta^*(s)}{\theta(s)},
\end{equation}

Eq.~\eqref{eq:axisymmetric_integral_momentum} can be expanded into an
evolution equation for the momentum thickness:

\begin{equation}
    \boxed{
    \frac{d\theta}{ds}
    =
    \frac{C_f}{2}
    -
    \theta
    \left[
        \frac{d\ln P_w}{ds}
        +
        \frac{d\ln\rho_e}{ds}
        +
        (2+H)\frac{d\ln U_e}{ds}
    \right]
    }.
    \label{eq:momentum_thickness_evolution}
\end{equation}

Since \(P_w=2\pi R_w\),

\begin{equation}
    \frac{d\ln P_w}{ds}
    =
    \frac{d\ln R_w}{ds},
\end{equation}

and the equation may equivalently be written as

\begin{equation}
    \boxed{
    \frac{d\theta}{ds}
    =
    \frac{C_f}{2}
    -
    \theta
    \left[
        \frac{d\ln R_w}{ds}
        +
        \frac{d\ln\rho_e}{ds}
        +
        (2+H)\frac{d\ln U_e}{ds}
    \right]
    }.
    \label{eq:momentum_thickness_evolution_radius}
\end{equation}

The terms in this equation represent, respectively,

\begin{itemize}
    \item production of momentum thickness by wall shear, \(C_f/2\);
    \item the effect of the changing axisymmetric wetted perimeter;
    \item the effect of the streamwise variation in edge density;
    \item the effect of core-flow acceleration and pressure gradient.
\end{itemize}

The integral model therefore transports the one-dimensional state
\(\theta(s)\), provided that closure relations are supplied for \(C_f(s)\) and
\(H(s)\). The profile closure supplies \(H(s)\), while the wall-shear closure
supplies \(C_f(s)\).

After integrating \(\theta(s)\), the displacement and representative physical
thicknesses are recovered from

\begin{equation}
    \boxed{
    \delta^*(s)=H(s)\theta(s)
    }
\end{equation}

and

\begin{equation}
    \boxed{
    \delta(s)=\frac{\theta(s)}{I_\theta(s)}
    }.
\end{equation}

\subsection{Skin Friction Closure Problem}

As it can be seen, skin friction can be calculated as: 

\begin{equation}
    \boxed{
    C_f(s)
    =
    \frac{2\tau_w(s)}
         {\rho_e(s)U_e^2(s)}
    }
    \label{eq:skin_friction_coefficient}
\end{equation}

Equation~\eqref{eq:skin_friction_closure_problem} is a definition rather than a
closure: evaluating it still requires \(\tau_w\). The Quick mode estimates wall
shear with an Eckert-reference-temperature turbulent flat-plate correlation and
the assumed one-seventh-power profile. This deliberately weak and inexpensive
closure is used only for screening. The recommended final calculation obtains
\(\tau_w\) from the molecular wall gradient resolved by the BLIMP-inspired model
described in Section~\ref{sec:blimp_cebeci_smith}.
